\section{Sheaves of \(k\)-modules: sheafification, Ext, and the structure sheaf as a sheaf of \(k\)-modules}
\label{mk-chap:Cohomology_StructureSheafModuleK}

Let \(C\) be a scheme over \(\Spec k\). The structure morphism makes every ring of sections \(\Gamma(C,U)\) a \(k\)-algebra, compatibly with restriction. Thus \(\mathcal O_C\) is a sheaf of \(k\)-modules. Since the corresponding sheaf category is \(k\)-linear and abelian, its \(\Ext\) groups carry canonical \(k\)-module structures.

Sections \crefrange{mk-sec:moduleK_prereqs}{mk-sec:moduleK_ext} establish sheafification and \(\Ext\) for sheaves of \(k\)-modules. Section \ref{mk-sec:scheme_toModuleKSheaf} constructs the \(k\)-module-valued structure sheaf from the algebra maps induced by the structure morphism.

\section{Sheafification of \(k\)-module-valued sheaves on opens}
\label{mk-sec:moduleK_prereqs}

\begin{theorem}[Sheafification on the topology of opens, \(k\)-module valued]\label{mk-thm:HasSheafify_Opens_ModuleCatK}
  \dcref{custom}

  Let \(k\) be a commutative ring and \(X\) a topological space. Then the Grothendieck topology of opens of \(X\) admits sheafification for sheaves valued in the category of \(k\)-modules.
\end{theorem}

\begin{proof}


The category of \(k\)-modules is Grothendieck abelian. Therefore module-valued presheaves on the site of opens admit the usual sheafification, whose functor is left exact.
\end{proof}

\section{\(\Ext\) for \(k\)-module-valued sheaves on opens}
\label{mk-sec:moduleK_ext}

\begin{theorem}[\(\Ext\) on the topology of opens, \(k\)-module valued]\label{mk-thm:HasExt_Sheaf_Opens_ModuleCatK}
  \dcref{custom}

  \uses{mk-thm:HasSheafify_Opens_ModuleCatK}

  Let \(k\) be a commutative ring and \(X\) a topological space. Then the category of sheaves of \(k\)-modules on the topology of opens of \(X\) admits \(\Ext\) groups.
\end{theorem}

\begin{proof}


By \ref{mk-thm:HasSheafify_Opens_ModuleCatK}, sheaves of \(k\)-modules form a Grothendieck abelian category. It has enough injectives, so the right derived functors of \(\Hom\) define \(\Ext\).
\end{proof}

\section{The structure sheaf of a \(\Spec k\)-scheme as a sheaf of \(k\)-modules}
\label{mk-sec:scheme_toModuleKSheaf}

For \(C\) a scheme over \(\Spec k\), the structure morphism \(C \to \Spec k\) induces an algebra map \(k \to \Gamma(C, U)\) for every open \(U \subseteq C\). These maps equip the structure sheaf with the structure of a sheaf of \(k\)-modules.

\subsection{The algebra map from the structure morphism}

For each open \(U \subseteq C\), the structure morphism induces a chain of ring maps
\[
  k \xrightarrow{\cong} \Gamma(\Spec k, \top) \xrightarrow{C.\mathrm{hom}} \Gamma(C, C.\mathrm{hom}^{-1}\top) \xrightarrow{\rho} \Gamma(C, U),
\]
where \(\rho\) is the restriction along \(U \le C.\mathrm{hom}^{-1}\top = \top\).

\begin{definition}[Algebra map on a section]\label{mk-def:Scheme_kToSection}
  \dcref{custom}


  Let \(k\) be a commutative ring and \(C\) a scheme over \(\Spec k\). For any open \(U\) of \(C\), define the ring map \(k \to \Gamma(C, U)\) to be the composition of the inverse of the canonical isomorphism \(\Gamma(\Spec k, \top) \cong k\), the structure morphism map on global sections, and the restriction from \(\top\) to \(U\).
\end{definition}

\begin{definition}[Algebra structure on a section]\label{mk-def:Scheme_algebraSection}
  \dcref{custom}

  \uses{mk-def:Scheme_kToSection}

  The ring map of \ref{mk-def:Scheme_kToSection} exhibits \(\Gamma(C, U)\) as a \(k\)-algebra.
\end{definition}

\begin{lemma}[Defining identity for the algebra map]\label{mk-lem:Scheme_algebraMap_eq_kToSection}
  \dcref{custom}

  \uses{mk-def:Scheme_algebraSection, mk-def:Scheme_kToSection}

  Under the algebra structure of \ref{mk-def:Scheme_algebraSection}, the structure map \(k \to \Gamma(C, U)\) equals the ring map of \ref{mk-def:Scheme_kToSection}.
\end{lemma}

\begin{proof}

  The algebra structure of \ref{mk-def:Scheme_algebraSection} is defined through the ring homomorphism of \ref{mk-def:Scheme_kToSection}; its structure map is therefore that homomorphism.
\end{proof}

\subsection{Naturality of the algebra map}

\begin{lemma}[Restriction is compatible with the algebra map]\label{mk-lem:Scheme_kToSection_naturality}
  \dcref{custom}

  \uses{mk-def:Scheme_kToSection}

  For any inclusion \(V \le U\) of opens in \(C\), the algebra map commutes with restriction: the composition \(k \to \Gamma(C, U) \to \Gamma(C, V)\) equals the direct map \(k \to \Gamma(C, V)\).
\end{lemma}

\begin{proof}

  Both sides factor through the structure morphism map on global sections. The remaining triangle commutes by functoriality of the presheaf on composable inclusions.
\end{proof}

\begin{lemma}[Algebra-map naturality]\label{mk-lem:Scheme_algebraMap_naturality}
  \dcref{custom}

  \uses{mk-lem:Scheme_algebraMap_eq_kToSection, mk-lem:Scheme_kToSection_naturality}

  For any inclusion \(V \le U\) of opens in \(C\) and any \(r \in k\), the restriction map sends the image of \(r\) in \(\Gamma(C, U)\) to the image of \(r\) in \(\Gamma(C, V)\).
\end{lemma}

\begin{proof}

  Replace the two algebra structure maps by the ring maps of \ref{mk-def:Scheme_kToSection} using \ref{mk-lem:Scheme_algebraMap_eq_kToSection}. Their compatibility with restriction is exactly \ref{mk-lem:Scheme_kToSection_naturality}.
\end{proof}

\subsection{The presheaf of \(k\)-modules and its sheaf condition}

\begin{definition}[Presheaf of \(k\)-modules from \(\mathcal O_C\)]\label{mk-def:Scheme_toModuleKPresheaf}
  \dcref{custom}

  \uses{mk-def:Scheme_algebraSection, mk-lem:Scheme_algebraMap_naturality}

  Define the presheaf of \(k\)-modules on \(C\) by sending each open \(U\) to the \(k\)-module \(\Gamma(C, U)\), with the \(k\)-action given by the algebra structure of \ref{mk-def:Scheme_algebraSection}, and on an inclusion \(V \le U\) by the restriction map viewed as a \(k\)-linear map. The identity and composition laws are those of the structure sheaf.
\end{definition}

\begin{theorem}[Sheaf condition]\label{mk-thm:Scheme_toModuleKPresheaf_isSheaf}
  \dcref{custom}

  \uses{mk-def:Scheme_toModuleKPresheaf}

  The presheaf of \ref{mk-def:Scheme_toModuleKPresheaf} is a sheaf for the Grothendieck topology of opens of \(C\).
\end{theorem}

\begin{proof}

  The forgetful functor from \(k\)-modules to sets creates limits and reflects the sheaf condition. The underlying presheaf of sets is the underlying presheaf of the structure sheaf, which is a sheaf because \(\mathcal O_C\) is a sheaf of commutative rings.
\end{proof}

\subsection{The structure sheaf as a sheaf of \(k\)-modules}

\begin{definition}[Structure sheaf as a sheaf of \(k\)-modules]\label{mk-def:Scheme_toModuleKSheaf}
  \dcref{custom}

  \uses{mk-def:Scheme_toModuleKPresheaf, mk-thm:Scheme_toModuleKPresheaf_isSheaf}

  Bundle \ref{mk-def:Scheme_toModuleKPresheaf} and \ref{mk-thm:Scheme_toModuleKPresheaf_isSheaf} into the sheaf of \(k\)-modules on \(C\).
\end{definition}

\subsection{Forget-and-recover: bridging \(k\)-module and abelian-group values}

\begin{definition}[Forget-and-recover natural iso for the \(k\)-module structure sheaf]\label{mk-def:Scheme_toModuleKSheaf_forgetCompare}
  \dcref{custom}

  \uses{mk-def:Scheme_toModuleKSheaf, mk-def:Scheme_toAbSheaf}

  Let \(C\) be a scheme over \(\Spec k\). Then forgetting the \(k\)-module structure from the sheaf of \ref{mk-def:Scheme_toModuleKSheaf} yields a sheaf of abelian groups canonically isomorphic to the abelian-group-valued structure sheaf of \ref{mk-def:Scheme_toAbSheaf}.
\end{definition}

\begin{proof}

  At every open \(U\), both sheaves have the abelian group underlying \(\Gamma(C,U)\), and both restriction morphisms are the underlying additive maps of the usual restrictions. These identifications commute with restriction and hence define the claimed natural isomorphism.
\end{proof}

\begin{lemma}[Object-evaluation lemma for the presheaf of \(k\)-modules]\label{mk-lem:Scheme_toModuleKPresheaf_obj}
  \dcref{custom}

  \uses{mk-def:Scheme_toModuleKPresheaf}

  For every open \(U\) of \(C\), the value of the presheaf of \ref{mk-def:Scheme_toModuleKPresheaf} at \(U\) is the \(k\)-module with carrier \(\Gamma(C, U)\).
\end{lemma}

\begin{proof}

  This is precisely the object assignment in \ref{mk-def:Scheme_toModuleKPresheaf}.
\end{proof}

\section{Sheaf cohomology with \(k\)-module coefficients}
\label{mk-sec:HModule}

In a \(k\)-linear abelian category, \(\Ext\) groups carry a canonical \(k\)-module structure. Applying this to sheaves of \(k\)-modules gives \(k\)-linear sheaf cohomology.

\begin{definition}[Sheaf cohomology with \(k\)-module coefficients]\label{mk-def:Scheme_HModule}
  \dcref{custom}

  \uses{mk-thm:HasExt_Sheaf_Opens_ModuleCatK, mk-thm:HasSheafify_Opens_ModuleCatK}

  Let \(k\) be a field, \(C\) a Grothendieck site admitting sheafification of \(k\)-modules and \(\Ext\). For a sheaf \(F\) of \(k\)-modules and \(n \in \mathbb N\), define the \(k\)-module-valued sheaf cohomology
  \[
    H^n_{\Module}(F) \;:=\; \Ext^n\bigl(\text{constant sheaf at }k,\; F\bigr).
  \]
  This automatically carries a \(k\)-module structure.
\end{definition}

\section{The \(H^0\) algebraic bridge}
\label{mk-sec:HModule_zero}

After \ref{mk-def:Scheme_HModule} the cohomology groups \(H^n_{\Module}(F)\) are well-defined and carry the canonical \(k\)-module structure. The next algebraic preliminary is the identification of the degree-zero piece \(H^0_{\Module}(F)\) with a \(\Hom\) group: the bridge from sheaf cohomology to representable Hom-functors.

\begin{definition}[\(H^0\) as a \(k\)-linear Hom group]\label{mk-def:Scheme_HModule_zero_linearEquiv}
  \dcref{custom}

  \uses{mk-def:Scheme_HModule}

  Let \(k\) be a field, \(C\) a Grothendieck site admitting sheafification of \(k\)-modules and \(\Ext\), and \(F\) a sheaf of \(k\)-modules. Then there is a canonical \(k\)-linear equivalence
  \[
    H^0_{\Module}(F) \;\simeq_k\; \Hom(\text{constant sheaf at }k,\; F).
  \]
\end{definition}

\begin{proof}

  In any abelian category, \(\Ext^0(A,F)\) is canonically \(\Hom(A,F)\). This identification respects scalar multiplication in a \(k\)-linear category; take \(A\) to be the constant sheaf at \(k\).
\end{proof}

\section{\v{C}ech cohomology of the structure sheaf}
\label{mk-sec:moduleK_cech}

The \v{C}ech complex of a sheaf with respect to an indexed open cover is the standard alternating-coface complex; its cohomology in each degree is a \(k\)-module.

\begin{definition}[\v{C}ech complex of the structure sheaf]\label{mk-def:Scheme_cechCochain_OC}
  \dcref{custom}

  \uses{mk-def:Scheme_toModuleKSheaf}
  Let \(C\) be a scheme over \(\Spec k\) and \(\mathcal U\) an indexed family of opens of \(C\). Define the \v{C}ech cochain complex of the structure sheaf with respect to \(\mathcal U\) to be the standard \v{C}ech complex applied to the underlying presheaf of the \(k\)-module-valued structure sheaf. This is a cochain complex valued in \(k\)-modules, indexed by \(\mathbb N\).
\end{definition}

\begin{definition}[\v{C}ech cohomology of the structure sheaf]\label{mk-def:Scheme_cechCohomology_OC}
  \dcref{custom}

  \uses{mk-def:Scheme_cechCochain_OC}
  With \((C, \mathcal U)\) as in \ref{mk-def:Scheme_cechCochain_OC}, define the \(n\)-th \v{C}ech cohomology to be the \(n\)-th homology of the \v{C}ech cochain complex. The result lives in the category of \(k\)-modules.
\end{definition}

\section{Cohomology of an open in \(k\)-module flavor}
\label{mk-sec:hmodule_prime}

For an object \(X\) of the site, cohomology over \(X\) is represented by the sheafification of the free \(k\)-module on the representable presheaf at \(X\).

\begin{definition}[Cohomology of an open with \(k\)-module coefficients]\label{mk-def:Scheme_HModule_prime}
  \dcref{custom}

  \uses{mk-thm:HasExt_Sheaf_Opens_ModuleCatK, mk-thm:HasSheafify_Opens_ModuleCatK}

  Let \(k\) be a field, \(C\) a Grothendieck site admitting sheafification of \(k\)-modules and \(\Ext\), and \(F\) a sheaf of \(k\)-modules. For \(n \in \mathbb N\) and \(X\) an object of \(C\):
  \[
    H^{\prime n}(X, F) \;:=\; \Ext^n\bigl(\text{sheafified representable at }X,\; F\bigr) \;\in\; \mathrm{Type}.
  \]
  The underlying type carries a \(k\)-module structure through the same \(\Ext\) mechanism that powers \(H_{\Module}\).
\end{definition}

\section{The \(H^0\) algebraic bridge for cohomology of an open}
\label{mk-sec:hmodule_prime_zero}

Parallel to \ref{mk-def:Scheme_HModule_zero_linearEquiv}, we identify the degree-zero piece of cohomology of an open with a Hom group.

\begin{definition}[\(H^0\) as a \(k\)-linear Hom group, open flavor]\label{mk-def:Scheme_HModule_prime_zero_linearEquiv}
  \dcref{custom}

  \uses{mk-def:Scheme_HModule_prime}

  Let \(k\) be a field, \(C\) a Grothendieck site, \(X\) an object of \(C\), and \(F\) a sheaf of \(k\)-modules. Then there is a canonical \(k\)-linear equivalence
  \[
    H^{\prime 0}(X, F) \;\simeq_k\; \Hom(\text{sheafified representable at }X,\; F).
  \]
\end{definition}

\begin{proof}

  Apply the canonical \(k\)-linear identification \(\Ext^0(A,F) \cong \Hom(A,F)\) with \(A\) the sheafified representable at \(X\).
\end{proof}

Together with \ref{mk-def:Scheme_HModule_zero_linearEquiv}, this identifies degree-zero global and open-local cohomology with their corresponding representable Hom functors.


After the \(H^0\) bridges identify cohomology with Hom groups, the immediate companion is transport of finite-dimensionality: if the Hom group is finite-dimensional over \(k\), so is the cohomology group.

\begin{theorem}[\(k\)-finite transport for global \(H^0\)]\label{mk-thm:Scheme_module_finite_HModule_zero}
  \dcref{custom}

  \uses{mk-def:Scheme_HModule_zero_linearEquiv}

  Let \(k\) be a field, \(C\) a Grothendieck site admitting sheafification of \(k\)-modules and \(\Ext\), and \(F\) a sheaf of \(k\)-modules. If the \(k\)-linear Hom group \(\Hom(\text{constant sheaf at }k, F)\) is finite-dimensional over \(k\), then \(H^0_{\Module}(F)\) is also finite-dimensional over \(k\).
\end{theorem}

\begin{proof}

  Transport finite-dimensionality across the symmetric form of the \(H^0\) bridge.
\end{proof}

\begin{theorem}[\(k\)-finite transport for open \(H^0\)]\label{mk-thm:Scheme_module_finite_HModule_prime_zero}
  \dcref{custom}

  \uses{mk-def:Scheme_HModule_prime_zero_linearEquiv}
  Let \(k\) be a field, \(C\) a Grothendieck site, \(X\) an object of \(C\), and \(F\) a sheaf of \(k\)-modules. If the \(k\)-linear Hom group \(\Hom(\text{sheafified representable at }X, F)\) is finite-dimensional over \(k\), then \(H^{\prime 0}(X, F)\) is also finite-dimensional over \(k\).
\end{theorem}

\begin{proof}

  Transport finite-dimensionality across the symmetric form of the open \(H^0\) bridge.
\end{proof}


\begin{theorem}[Curve specialisation of global \(H^0\) finite-dimensionality]\label{mk-thm:Scheme_module_finite_HModule_zero_curve}
  \dcref{custom}

  \uses{mk-thm:Scheme_module_finite_HModule_zero, mk-def:Scheme_toModuleKSheaf}
  Let \(k\) be a field and \(C\) a scheme over \(\Spec k\). If the \(k\)-linear Hom group from the constant sheaf at \(k\) to the \(k\)-module-valued structure sheaf is finite-dimensional over \(k\), then \(H^0_{\Module}(\mathcal O_C)\) is finite-dimensional over \(k\).
\end{theorem}

\begin{theorem}[Curve specialisation of open \(H^0\) finite-dimensionality]\label{mk-thm:Scheme_module_finite_HModule_prime_zero_curve}
  \dcref{custom}

  \uses{mk-thm:Scheme_module_finite_HModule_prime_zero, mk-def:Scheme_toModuleKSheaf}
  Let \(k\) be a field, \(C\) a scheme over \(\Spec k\), and \(U\) an open of \(C\). If the \(k\)-linear Hom group from the sheafified representable at \(U\) to the \(k\)-module-valued structure sheaf is finite-dimensional over \(k\), then \(H^{\prime 0}(U, \mathcal O_C)\) is finite-dimensional over \(k\).
\end{theorem}


For a Mayer--Vietoris argument on a two-affine cover, one uses Serre's theorem that the positive-degree cohomology of a quasicoherent sheaf on each affine piece vanishes. We isolate the corresponding vanishing property for a sheaf of \(k\)-modules.

\begin{definition}[Affine cohomology vanishing]\label{mk-thm:Scheme_IsAffineHModuleVanishing}
  \dcref{custom}

  \uses{mk-def:Scheme_HModule_prime}
  Let \(k\) be a field, \(C\) a scheme over \(\Spec k\), and \(F\) a sheaf of \(k\)-modules on \(C\). We say that \(F\) has affine cohomology vanishing if, for every affine open \(U\) of \(C\) and every degree \(i > 0\), the group \(H^{\prime i}(U, F)\) is zero.
\end{definition}

\begin{theorem}[Finite-dimensionality from affine vanishing]\label{mk-thm:Scheme_module_finite_HModule_prime_of_isAffineHModuleVanishing}
  \dcref{custom}

  \uses{mk-thm:Scheme_IsAffineHModuleVanishing, mk-def:Scheme_HModule_prime}
  Let \(k\) be a field, \(C\) a scheme over \(\Spec k\), and \(F\) a sheaf of \(k\)-modules. Assume affine cohomology vanishing holds for \(F\). For any affine open \(U\) of \(C\) and any \(i > 0\), the cohomology group \(H^{\prime i}(U, F)\) is finite-dimensional over \(k\).
\end{theorem}

\begin{proof}

  A subsingleton \(k\)-module is generated by the empty set, hence finitely generated and finitely presented, hence finite-dimensional.
\end{proof}


Global sections of a proper scheme can be finite-dimensional even though sections on its proper affine opens are not; for example, \(\Gamma(\mathbb A^1_k,\mathcal O)=k[t]\). Accordingly, the relevant degree-zero finiteness condition concerns only the global Hom group.

\begin{definition}[Global \(H^0\) Hom-finiteness]\label{mk-thm:Scheme_IsHModuleHomFinite}
  \dcref{custom}


  Let \(k\) be a field, \(C\) a scheme over \(\Spec k\), and \(F\) a sheaf of \(k\)-modules on \(C\). We say that \(F\) has global \(H^0\) Hom-finiteness if \(\Hom(\text{constant sheaf at }k, F)\) is finite-dimensional over \(k\).
\end{definition}

\begin{theorem}[Finite-dimensionality of global \(H^0\) from global Hom-finiteness]\label{mk-thm:Scheme_module_finite_HModule_zero_of_isHModuleHomFinite}
  \dcref{custom}

  \uses{mk-thm:Scheme_IsHModuleHomFinite, mk-thm:Scheme_module_finite_HModule_zero}
  Let \(k\) be a field, \(C\) a scheme over \(\Spec k\), and \(F\) a sheaf of \(k\)-modules. Assume global \(H^0\) Hom-finiteness holds for \(F\). Then \(H^0_{\Module}(F)\) is finite-dimensional over \(k\).
\end{theorem}

\begin{theorem}[Curve specialisation]\label{mk-thm:Scheme_module_finite_HModule_zero_of_isHModuleHomFinite_curve}
  \dcref{custom}

  \uses{mk-thm:Scheme_module_finite_HModule_zero_of_isHModuleHomFinite, mk-def:Scheme_toModuleKSheaf}
  Let \(k\) be a field and \(C\) a scheme over \(\Spec k\). If the global morphism group from the constant sheaf at \(k\) to \(\mathcal O_C\), viewed as a sheaf of \(k\)-modules, is finite-dimensional, then \(H^0_{\Module}(\mathcal O_C)\) is finite-dimensional over \(k\).
\end{theorem}

\begin{proof}
  Apply \ref{mk-thm:Scheme_module_finite_HModule_zero_of_isHModuleHomFinite} to the \(k\)-module-valued structure sheaf.
\end{proof}

\begin{remark}
  On a proper geometrically connected curve, Stein factorization gives \(\Gamma(C, \mathcal O_C) = k\), which is finite-dimensional over \(k\).
\end{remark}


The geometric input for global Hom-finiteness is Stein's theorem: on a proper integral scheme over a field, global sections are finite-dimensional.

\begin{theorem}[Module-finiteness of the top-section map of a universally closed morphism]
  \label{mk-thm:finite_appTop_of_universallyClosed}
  \dcref{custom}


  Let \(K\) be a commutative ring and \(X\) an integral scheme equipped with a
  universally closed, locally-of-finite-type morphism \(f : X \to \Spec K\).
  Then the ring map on top sections \(f^\sharp_{\top}\) is module-finite.

\end{theorem}

\begin{theorem}[Stein finiteness of global sections]\label{mk-thm:Scheme_module_finite_globalSections_of_isProper}
  \dcref{custom}

  \uses{mk-thm:finite_appTop_of_universallyClosed, mk-def:Scheme_kToSection}

  Let \(k\) be a field and \(C\) an integral scheme over \(\Spec k\) with proper structure morphism. Then \(\Gamma(C, \mathcal O_C)\) is finite-dimensional over \(k\).
\end{theorem}

\begin{proof}

  Apply the theorem that for an integral scheme with a universally closed, locally finite-type morphism to \(\Spec k\), the structure-morphism ring map on global sections is module-finite. Transfer the base ring from \(\Gamma(\Spec k, \top)\) to \(k\) along the canonical ring isomorphism.
\end{proof}

\section{Global-sections evaluation isomorphism}
\label{mk-sec:gammaObj_linearEquiv_top}

To bridge from global sections (a concrete \(k\)-module) to the abstract Hom-from-constant-sheaf form, we lift the natural isomorphism between global sections and evaluation at the terminal object to a \(k\)-linear equivalence.

\begin{theorem}[Global-sections-to-top linear equivalence]\label{mk-thm:Scheme_SheafGammaObj_linearEquiv_top}
  \dcref{custom}


  Let \(k\) be a field, \(X\) a topological space, and \(F\) a sheaf of \(k\)-modules on \(X\). Then there is a \(k\)-linear equivalence between the global-sections object of \(F\) and the evaluation of \(F\) at the top open \(\top\).
\end{theorem}

\begin{theorem}[Finite-dimensionality on global-sections object]\label{mk-thm:Scheme_module_finite_gammaObj_of_isProper}
  \dcref{custom}

  \uses{mk-thm:Scheme_module_finite_globalSections_of_isProper, mk-thm:Scheme_SheafGammaObj_linearEquiv_top}
  Let \(k\) be a field and \(C\) an integral scheme over \(\Spec k\) with proper structure morphism. Then the global-sections object of the \(k\)-module-valued structure sheaf is finite-dimensional over \(k\).
\end{theorem}

\begin{proof}

  Combine \ref{mk-thm:Scheme_module_finite_globalSections_of_isProper} with the linear equivalence of \ref{mk-thm:Scheme_SheafGammaObj_linearEquiv_top}, transporting finite-dimensionality across the symmetric form.
\end{proof}

\section{Global Hom-finiteness for the structure sheaf}
\label{mk-sec:IsHModuleHomFinite_producer}

The constant-sheaf/global-sections adjunction and evaluation at \(1\in k\) identify the global Hom group with global sections. Stein finiteness then proves global \(H^0\) Hom-finiteness for the structure sheaf of a proper integral scheme.

\begin{lemma}[Additivity of the constant functor]\label{mk-lem:Functor_const_additive}
  \dcref{custom}

  Let \(C\) be a category and \(D\) a preadditive category. Then the constant functor
  \(\Functor.const\, C : D \to (C \to D)\) is additive.
\end{lemma}

\begin{proof}

For morphisms \(f,g\) in \(D\), the components of the constant natural transformation associated to \(f+g\) are all \(f+g\), which are also the components of the sum of the constant transformations associated to \(f\) and \(g\). The same argument applies to zero, so the constant functor is additive.
\end{proof}

\begin{lemma}[\(R\)-linearity of the constant functor]\label{mk-lem:Functor_const_linear}
  \dcref{custom}

  \uses{mk-lem:Functor_const_additive}
  Let \(C\) be a category, \(R\) a semiring, and \(D\) an \(R\)-linear preadditive category.
  Then the constant functor \(\Functor.const\, C : D \to (C \to D)\) is \(R\)-linear.
\end{lemma}

\begin{proof}

For \(r\in R\) and a morphism \(f\) of \(D\), both the constant transformation associated to \(r\cdot f\) and \(r\) times the constant transformation associated to \(f\) have component \(r\cdot f\) at every object. Equality of natural transformations is componentwise, so the constant functor is \(R\)-linear.
\end{proof}

\begin{lemma}[Linear lifting of adjoint functors]\label{mk-lemma:Adjunction_left_adjoint_linear}
  \dcref{custom}

  Let \(F \dashv G\) be an adjunction between \(R\)-linear preadditive categories. If \(G\) is \(R\)-linear, then \(F\) is \(R\)-linear.
\end{lemma}

\begin{proof}
  For a morphism \(f:X\to X'\), the adjunction sends \(h:FX\to FX'\) to \(G(h)\circ\eta_X\). Naturality of the unit identifies the images of \(F(r\cdot f)\) and \(r\cdot F(f)\), because \(G\) preserves scalar multiplication. Since the adjunction map is injective, these morphisms agree.
\end{proof}

\begin{lemma}[Right-adjoint version]\label{mk-lemma:Adjunction_right_adjoint_linear}
  \dcref{custom}

  Dual of \ref{mk-lemma:Adjunction_left_adjoint_linear}: if \(F\) is \(R\)-linear, then \(G\) is \(R\)-linear.
\end{lemma}

\begin{proof}
  For \(f:Y\to Y'\), the inverse adjunction map sends \(h:GY\to GY'\) to \(\varepsilon_{Y'}\circ F(h)\). Naturality of the counit identifies the images of \(G(r\cdot f)\) and \(r\cdot G(f)\), because \(F\) preserves scalar multiplication. Injectivity of the inverse adjunction map gives the result.
\end{proof}

\begin{theorem}[Linear hom-equivalence from an adjunction]\label{mk-def:Adjunction_homLinearEquiv}
  \dcref{custom}


  Let \(F \dashv G\) be an adjunction between \(R\)-linear preadditive categories with \(G\) additive and \(R\)-linear. Then for all objects \(X\) and \(Y\), the additive equivalence of the adjunction upgrades to an \(R\)-linear equivalence between \(\Hom(FX, Y)\) and \(\Hom(X, GY)\).
\end{theorem}

\begin{proof}
  The adjunction sends \(f:FX\to Y\) to \(G(f)\circ\eta_X\). Preadditivity and the \(R\)-linearity of \(G\) show that this map preserves addition and scalar multiplication. The unit argument shows that \(F\) is \(R\)-linear, so the inverse map \(g\mapsto\varepsilon_Y\circ F(g)\) is also linear. Thus the adjunction bijection is an \(R\)-linear equivalence.
\end{proof}

\begin{theorem}[Constant-sheaf / global-sections linear equivalence]\label{mk-def:Scheme_constantSheafGammaHom_linearEquiv}
  \dcref{custom}

  \uses{mk-lemma:Adjunction_left_adjoint_linear, mk-lemma:Adjunction_right_adjoint_linear, mk-def:Adjunction_homLinearEquiv, mk-lem:Functor_const_linear}
  Let \(k\) be a field, \(C\) a Grothendieck site admitting sheafification of \(k\)-modules, and \(F\) a sheaf of \(k\)-modules. For any \(k\)-module \(X\), there is a \(k\)-linear equivalence between morphisms from the constant sheaf at \(X\) to \(F\), and \(k\)-linear maps from \(X\) to the global sections of \(F\).
\end{theorem}

\begin{proof}
  The constant-sheaf functor is left adjoint to global sections. The constant functor is \(k\)-linear by \ref{mk-lem:Functor_const_linear}, and linearity passes across the adjunction by \ref{mk-lemma:Adjunction_left_adjoint_linear} and \ref{mk-lemma:Adjunction_right_adjoint_linear}. The adjunction bijection is therefore \(k\)-linear by \ref{mk-def:Adjunction_homLinearEquiv}.
\end{proof}

\begin{theorem}[Hom-from-\(k\) evaluation]\label{mk-def:Scheme_homFromOne_linearEquiv}
  \dcref{custom}

  For any field \(k\) and \(k\)-module \(M\), there is a canonical \(k\)-linear equivalence between \(\Hom_k(k, M)\) and \(M\) given by evaluation at \(1\).
\end{theorem}

\begin{proof}
  Evaluation sends \(f\) to \(f(1)\). Its inverse sends \(m\in M\) to the linear map \(r\mapsto r\cdot m\). Linearity and the two inverse identities follow from \(f(r)=r\cdot f(1)\).
\end{proof}

\begin{theorem}[Global Hom-finiteness of the structure sheaf]\label{mk-inst:Scheme_instIsHModuleHomFinite_toModuleKSheaf}
  \dcref{custom}

  \uses{mk-thm:Scheme_module_finite_gammaObj_of_isProper, mk-def:Scheme_constantSheafGammaHom_linearEquiv, mk-def:Scheme_homFromOne_linearEquiv, mk-thm:Scheme_IsHModuleHomFinite}
  Let \(k\) be a field and \(C\) an integral scheme over \(\Spec k\) with proper structure morphism. Then global \(H^0\) Hom-finiteness holds for the \(k\)-module-valued structure sheaf of \(C\).
\end{theorem}

\begin{proof}

  Compose the constant-sheaf/global-sections linear equivalence (specialised to \(X = k\)) with the Hom-from-\(k\) evaluation to obtain a linear equivalence between the global Hom group and the global-sections object. \ref{mk-thm:Scheme_module_finite_gammaObj_of_isProper} provides finite-dimensionality on the latter; transport across the symmetric form of the equivalence.
\end{proof}

\section{\v{C}ech complexes for arbitrary sheaves}
\label{mk-sec:moduleK_cech_parameterised}

The \v{C}ech construction applies to an arbitrary sheaf of \(k\)-modules; the structure-sheaf construction is its specialization at \(F=\mathcal O_C\).

\begin{definition}[Parameterised \v{C}ech complex]\label{mk-def:Scheme_cechCochain}
  \dcref{custom}


  Let \(C\) be a scheme over \(\Spec k\), \(F\) a sheaf of \(k\)-modules on \(C\), and \(\mathcal U\) an indexed family of opens. Define the \v{C}ech cochain complex of \(F\) with respect to \(\mathcal U\) to be the standard \v{C}ech complex applied to the underlying presheaf of \(F\). This is a cochain complex valued in \(k\)-modules.
\end{definition}

\begin{definition}[Parameterised \v{C}ech cohomology]\label{mk-def:Scheme_cechCohomology}
  \dcref{custom}

  \uses{mk-def:Scheme_cechCochain}
  With \((C, F, \mathcal U)\) as in \ref{mk-def:Scheme_cechCochain}, define the \(n\)-th \v{C}ech cohomology to be the \(n\)-th homology of the \v{C}ech cochain complex.
\end{definition}

\begin{theorem}[\v{C}ech complex bridge to structure-sheaf specialisation]\label{mk-thm:Scheme_cechCochain_OC_eq}
  \dcref{custom}

  \uses{mk-def:Scheme_cechCochain_OC, mk-def:Scheme_cechCochain, mk-def:Scheme_toModuleKSheaf}
  The structure-sheaf \v{C}ech complex equals the parameterised \v{C}ech complex at \(F = \mathcal O_C\).
\end{theorem}

\begin{proof}
  Both sides are the standard \v{C}ech complex of the same \(k\)-module-valued structure sheaf with respect to the same indexed family of opens.
\end{proof}

\begin{theorem}[\v{C}ech cohomology bridge to structure-sheaf specialisation]\label{mk-thm:Scheme_cechCohomology_OC_eq}
  \dcref{custom}

  \uses{mk-def:Scheme_cechCohomology_OC, mk-def:Scheme_cechCohomology, mk-def:Scheme_toModuleKSheaf}
  The structure-sheaf \v{C}ech cohomology equals the parameterised \v{C}ech cohomology at \(F = \mathcal O_C\).
\end{theorem}

\begin{proof}
  Both sides are the degree-\(n\) cohomology of the same structure-sheaf \v{C}ech complex, so they agree.
\end{proof}

\section{\v{C}ech acyclicity}
\label{mk-sec:moduleK_cech_acyclic_cover}

Positive-degree vanishing of \v{C}ech cohomology transports to derived cohomology whenever the two theories are related by a comparison isomorphism.

\begin{definition}[\v{C}ech acyclicity]\label{mk-def:Scheme_IsCechAcyclicCover}
  \dcref{custom}


  Let \(k\) be a field, \(C\) a scheme over \(\Spec k\), \(F\) a sheaf of \(k\)-modules on \(C\), and \(\mathcal U\) an indexed family of opens. The cover \(\mathcal U\) is \v{C}ech-acyclic for \(F\) if the \v{C}ech cohomology vanishes in positive degrees: for every \(n \ge 1\), the \(n\)-th \v{C}ech cohomology group is a subsingleton.
\end{definition}

\begin{theorem}[Derived vanishing from \v{C}ech acyclicity]\label{mk-thm:Scheme_subsingleton_HModule_prime_supr_of_isCechAcyclicCover}
  \dcref{custom}


  Let \(k\) be a field, \(C\) a scheme over \(\Spec k\), \(F\) a sheaf of \(k\)-modules, and \(\mathcal U\) a \v{C}ech-acyclic cover for \(F\). Suppose there is a \(k\)-linear comparison isomorphism between \v{C}ech cohomology and derived cohomology of the cover supremum for every degree. Then for every \(n \ge 1\), the derived cohomology of the cover supremum is a subsingleton.
\end{theorem}

\begin{proof}

  Transport the subsingleton structure along the symmetric form of the comparison isomorphism.
\end{proof}

\begin{theorem}[Curve specialisation]\label{mk-thm:Scheme_subsingleton_HModule_prime_supr_of_isCechAcyclicCover_curve}
  \dcref{custom}


  Let \(C\) be a scheme over \(\Spec k\) and \(\mathcal U\) a \v{C}ech-acyclic cover for \(\mathcal O_C\). Suppose \v{C}ech and derived cohomology of the cover supremum are \(k\)-linearly isomorphic in every degree. Then, for every \(n\geq 1\),
  \[
    H^{\prime n}\!\left(\bigsqcup_i \mathcal U_i,\mathcal O_C\right)=0.
  \]
\end{theorem}

\begin{proof}
  The positive-degree \v{C}ech group is zero by acyclicity. Its image under the assumed linear isomorphism is therefore the zero derived-cohomology group of the cover supremum.
\end{proof}

\section{Cohomological consequences}

The constructions above give the following compatible descriptions of structure-sheaf cohomology:

\begin{itemize}
  \item For \(C\) a scheme over \(\Spec k\), the cohomology \(H^n_{\Module}(\mathcal O_C)\) carries a canonical \(k\)-module structure.
  \item Forgetting the \(k\)-module structure recovers the abelian-group cohomology of \ref{mk-chap:Cohomology_StructureSheafAb}; their underlying abelian groups agree.
  \item Finite-dimensionality can therefore be expressed as \(\dim_k H^n_{\Module}(\mathcal O_C)<\infty\), and in degree zero it is transported across the Hom and global-sections equivalences above.
\end{itemize}
